\begin{table}[t]
\centering
\caption{Ablation study on 10 representative sparse instances. Relative change is measured against LR-TA. Negative percentages indicate improvement.}
\label{tab:ablation}
\begin{tabular}{lrrr}
\toprule
Variant & Mean BW & Relative change & Mean RT (s) \\
\midrule
LR without add-back & 4525.1 & +5.94\% & 0.017 \\
LR-TA & 4271.5 & +0.00\% & 0.047 \\
WMSF seed & 4332.5 & +1.43\% & 0.040 \\
Best seed, no refinement & 4271.5 & +0.00\% & 0.069 \\
IPSNS, 50 iterations & 4239.2 & -0.76\% & 0.778 \\
IPSNS, 400 iterations (default) & 4239.2 & -0.76\% & 5.734 \\
IPSNS, no SCC priority & 4239.1 & -0.76\% & 5.727 \\
\bottomrule
\end{tabular}
\par\vspace{2pt}\parbox{0.94\linewidth}{\footnotesize The add-back phase accounts for the largest mean improvement (about 5.9\%), while IPSNS contributes a further 0.75\% on this subset. The equal 50- and 400-iteration means support the default iteration budget only as subset-level evidence, not as a universal tuning rule.}
\end{table}
